\section{Выводы}
Выполнив седьмую лабораторную работу по курсу \enquote{Дискретный анализ}, я познакомился с
методом динамического программирования и применил его к задаче о преобразовании числа в единицу с
минимальной суммарной стоимостью операций.

Главное, что я понял: сначала нужно проверить, что задача обладает оптимальной подструктурой и перекрывающимися
подзадачами, и только потом выбирать метод. Жадный выбор в этой задаче не работает, потому что локально выгодная
операция (например, деление) не всегда ведёт к минимальной общей стоимости. Решение через рекуррентное
соотношение $dp[i] = i + \min(dp[i-1], dp[i/2], dp[i/3])$ корректно, так как оптимальный путь для числа $i$ состоит 
з первой операции и оптимального пути для полученного числа.

Тест производительности показал, почему запоминание подзадач необходимо. Наивная рекурсия пересчитывает одни и те
же значения, и уже при $n = 750$ ей требуется около трёх секунд, тогда как динамическое программирование решает эту же
задачу за единицы микросекунд. При $n = 10^7$ ДП работает около 0{,}05 с, а время на один элемент таблицы остаётся в пределах
нескольких наносекунд, что согласуется с оценкой $O(n)$. Я также заметил, что на очень малых входах измерения неустойчивы из-за
погрешности таймера и накладных расходов, а на больших заметно влияет кэш процессора, поэтому асимптотику нужно проверять на
достаточно больших данных и с несколькими повторами замера.

Отдельно я научился восстанавливать не только значение оптимума, но и сам путь: для этого в каждой ячейке
таблицы достаточно хранить операцию, на которой достигнут минимум, и пройти по ним от $n$ до единицы. Для проверки
корректности я сравнил результаты ДП с наивным перебором на малых $n$ и пересчитал стоимость по восстановленной последовательности операций.

К практическим наблюдениям отнесу выбор между восходящим и нисходящим подходом. Рекурсия с запоминанием
при $n = 10^7$ имела бы глубину порядка $n$ и привела бы к переполнению стека, поэтому итеративное заполнение
таблицы здесь надёжнее. Недостатком выбранного решения является память $O(n)$: при максимальном $n$ таблица
занимает около 80~МБ, и для более жёстких ограничений пришлось бы искать более экономный способ.

Метод динамического программирования применим во многих задачах: поиск наибольшей общей
подпоследовательности, расстояние редактирования, задача о рюкзаке, кратчайшие пути в графах,
разбиение последовательностей. Полученный опыт (определить состояние, записать рекуррентное соотношение,
выбрать порядок вычисления, восстановить ответ) пригодится для решения таких задач.
\pagebreak